\section{Introduction}
\label{sec:introduction}

% WORDNET: manual creation cf. le2019inferring, hearst1992hyponyms; hyponym = is-a relationship cf. hearst1992hyponyms
Concept hierarchies form the backbone of many expert ontologies.
In linguistic and lexical-semantic ontologies, they organize hyponym relationships under a common top node~\cite{le2019inferring} and structure information into categories, thereby supporting rule formulation as well as the retrieval from and reuse of knowledge bases~\cite{cimiano_learning_2005,yang2012personalized}.
However, the manual construction of high-quality domain-specific ontologies is costly and typically reflects only a snapshot of evolving knowledge~\cite{li2018unsupervised}.
Although manually curated ontologies are often regarded as the gold standard, automatic generation is desirable to reduce time and cost~\cite{Aznag2014LeveragingFC,risch2020hierarchical}.

% semantic dictionaries (Roget's Thesaurus)
% lexical databases (WordNet, FrameNet)
% Proof of concept on German prefices sporleder2002galois
% Proof of concept on inflectional properties of German nouns petersen:04
% reports on domestic violence (report id=obj, phrase=attr) Elzinga:09
% 1. intro to automatic concept hierarchy generation
Automatic concept hierarchy generation provides a scalable alternative and enables coverage of domains not represented in existing ontologies.
% 2. already structured data sources
Structured linguistic resources, such as the semantic dictionary Roget's Thesaurus~\cite{Priss:2004,old:04,Priss:96}, lexical databases such as WordNet~\cite{Priss:96,priss:98,GalnPez2014DiscoveringNS} and FrameNet~\cite{val:05}, and morphological feature systems~\cite{sporleder:02,petersen:04}, can be represented as formal contexts.
% 3. unstructured data sources, e.g., hearst1992hyponyms with Hearst patterns: automatically extract hyponym–hypernym relations; cimiano_learning_2005 with a multi-step NLP pipeline: parser/POS, lemmatizer ...
When applied to unstructured data sources such as text corpora, formal contexts must first be constructed from textual information. 
Approaches range from pattern-based extraction methods~\cite{hearst1992hyponyms,le2019inferring} to multi-step \acl{nlp} pipelines that derive binary object--attribute contexts from texts~\cite{cimiano_learning_2005}, as well as approaches using phrases directly as attributes~\cite{Elzinga:09}.
 


\ac{fca}-based concept hierarchy generation can be viewed as a form of conceptual clustering~\cite{Cimiano2004ComparingCD} that yields intensional descriptions of abstract concepts through their intents~\cite{cimiano_learning_2005}. 
Since each concept is characterized by an explicit set of shared attributes, the resulting hierarchy provides transparent, attribute-based explanations of the relationships between concepts, making \ac{fca}-based hierarchies inherently interpretable compared to similarity-based clustering methods.

Despite their scalability, automatically generated hierarchies may conflict with established expert ontologies, potentially reducing trust in either source.
Ensuring consistency with hierarchical background knowledge therefore requires mechanisms that guide and constrain the generation process.
Existing approaches (Section~\ref{sec:related_work}) include human-in-the-loop strategies~\cite{yang2012personalized}, the use of \aclp{llm}~\cite{kommineni2024from,zhang2025teleclass}, constrained \ac{fca}~\cite{belohlavek2011selecting,belohlavek2009formal,belohlavek2006formal,hidayat_constrained_fca_2024,mao_constrained_fca_2017}, and hierarchical constraints in similarity-based settings~\cite{bade2006personalized,bade_hierarchical_2014}. 

% Besides presentation purposes, \ac{fca} has been employed to identify potentially missing thesaurus entries in a specific expert ontology~\cite{zheng2020lexical,hao2021leveraging}. 

To date, the systematic integration of hierarchical background knowledge into \ac{fca}-based concept hierarchy generation through \ac{mlb} constraints has received limited attention.
We address this gap by incorporating \ac{mlb} constraints into the \ac{fca} framework to guide the construction of concept hierarchies (Section~\ref{sec:method}),\footnote{The code is available at \url{github.com/KlaraGtknst/fca-constrained-clustering}} thereby extending prior work that incorporates pairwise constraints~\cite{belohlavek2011selecting}.
In addition, we present a case study (Section~\ref{sec:case_study}) in which a concept lattice generated under \ac{mlb} constraints is compared to a concept lattice derived from an independent knowledge source using corpus-specific statistical topic modeling. 
The comparison is conducted across different levels of coherence, giving rise to the \textit{consistency} and \textit{coherence} lattices.

%\textcolor{red}{NOTE: Their conclusion includes future work on using humans for ontological relation validation (Attr. Explor.???). They also compare generated ontologies with human-modeled ontologies via Precision, Recall, F-measure...}

% Usually you
%Since this approach relies solely on syntactic structure present in the text corpus, the available knowledge is very limited.
%Hierarchical constraints can be thought of a tool to model certain degrees of certitude with respect to the information's validity.